\documentclass[11pt]{article}
\usepackage[margin=1in]{geometry}
\usepackage{amsmath}
\usepackage{amssymb}
\usepackage{tcolorbox}

\title{Problem 3: von Neumann-Morgenstern Axioms \\ (Starting from Scratch)}
\author{ECO 310 - For People Who Skipped Lecture}
\date{February 23, 2026}

\begin{document}

\maketitle

\section{What Even Are Lotteries?}

\textbf{Lottery} = A gamble with random outcomes

\textbf{Notation:} $(p_a, p_b, p_c)$ means:
\begin{itemize}
\item $p_a$ = probability of getting apple
\item $p_b$ = probability of getting banana
\item $p_c$ = probability of getting cherry
\item These probabilities add to 1: $p_a + p_b + p_c = 1$
\end{itemize}

\textbf{Examples:}
\begin{itemize}
\item $(1, 0, 0)$ = apple for sure (100\% apple, 0\% banana, 0\% cherry)
\item $(0, 1/2, 1/2)$ = flip a coin, 50\% banana or 50\% cherry
\item $(1/3, 1/3, 1/3)$ = equal chance of each fruit
\end{itemize}

\section{What Is Expected Utility?}

If you assign utility numbers $u_a, u_b, u_c$ to apple, banana, cherry, then:

\begin{tcolorbox}[colback=blue!5!white,colframe=blue!75!black,title=Expected Utility Formula]
\textbf{Expected Utility of lottery} $(p_a, p_b, p_c)$ is:
$$EU = p_a \cdot u_a + p_b \cdot u_b + p_c \cdot u_c$$

Just the weighted average of utilities!
\end{tcolorbox}

\textbf{Example:}
\begin{itemize}
\item $u_a = 10$ (apple), $u_b = 5$ (banana), $u_c = 0$ (cherry)
\item Lottery $(0, 1/2, 1/2)$: $EU = 0(10) + (1/2)(5) + (1/2)(0) = 2.5$
\item Lottery $(1, 0, 0)$: $EU = 1(10) + 0(5) + 0(0) = 10$
\end{itemize}

Apple for sure (EU=10) is better than the 50-50 banana/cherry lottery (EU=2.5).

\section{What Are vNM Axioms?}

\begin{tcolorbox}[colback=yellow!10!white,colframe=orange!75!black,title=Simple Version]
The vNM axioms are rules about preferences over lotteries.

If someone follows ALL these rules, you can find utility numbers that make their preferences match expected utility.

If they violate ANY rule, expected utility doesn't work for them.
\end{tcolorbox}

\subsection{The 3 Main Axioms}

\textbf{A.1 - Completeness:} You can compare any two lotteries (prefer one or indifferent)

\textbf{A.2 - Transitivity:} If A $\succ$ B and B $\succ$ C, then A $\succ$ C (no cycles)

\textbf{A.3 - Independence:} If you're indifferent between A and B, mixing both with lottery C (same way) keeps you indifferent:
$$\text{If } A \sim B \text{, then } \alpha A + (1-\alpha) C \sim \alpha B + (1-\alpha) C$$

\section{Problem 3 Setup}

Dave's preferences:
\begin{itemize}
\item $p_1 = (1, 0, 0)$ [apple for sure] $\sim$ $p_2 = (0, 1/2, 1/2)$ [50\% banana, 50\% cherry]
\item $p_3 = (1/2, 1/2, 0)$ [50\% apple, 50\% banana] $\succ$ $p_4 = (0, 2/3, 1/3)$ [2/3 banana, 1/3 cherry]
\end{itemize}

\textbf{Goal:} Show this violates vNM axioms (can't find consistent utilities).

\section{Step-by-Step Solution}

\subsection{Step 1: Assume vNM Holds}

If Dave satisfies vNM axioms, there exist utilities $u_a, u_b, u_c$ such that:
\begin{itemize}
\item Dave prefers lottery X over Y $\iff$ $EU(X) > EU(Y)$
\item Dave is indifferent $\iff$ $EU(X) = EU(Y)$
\end{itemize}

\subsection{Step 2: Use the Indifference}

Dave says: $p_1 \sim p_2$

This means: $EU(p_1) = EU(p_2)$

Calculate:
\begin{align}
EU(p_1) &= 1 \cdot u_a + 0 \cdot u_b + 0 \cdot u_c = u_a \\
EU(p_2) &= 0 \cdot u_a + \frac{1}{2} \cdot u_b + \frac{1}{2} \cdot u_c = \frac{1}{2}u_b + \frac{1}{2}u_c
\end{align}

Setting them equal:
\begin{equation}
u_a = \frac{1}{2}u_b + \frac{1}{2}u_c
\end{equation}

\begin{tcolorbox}[colback=green!5!white,colframe=green!75!black,title=Key Equation 1]
$$u_a = \frac{1}{2}u_b + \frac{1}{2}u_c$$
\end{tcolorbox}

\subsection{Step 3: Use the Strict Preference}

Dave says: $p_3 \succ p_4$

This means: $EU(p_3) > EU(p_4)$

Calculate:
\begin{align}
EU(p_3) &= \frac{1}{2} \cdot u_a + \frac{1}{2} \cdot u_b + 0 \cdot u_c = \frac{1}{2}u_a + \frac{1}{2}u_b \\
EU(p_4) &= 0 \cdot u_a + \frac{2}{3} \cdot u_b + \frac{1}{3} \cdot u_c = \frac{2}{3}u_b + \frac{1}{3}u_c
\end{align}

The inequality:
\begin{equation}
\frac{1}{2}u_a + \frac{1}{2}u_b > \frac{2}{3}u_b + \frac{1}{3}u_c
\end{equation}

\begin{tcolorbox}[colback=green!5!white,colframe=green!75!black,title=Key Inequality]
$$\frac{1}{2}u_a + \frac{1}{2}u_b > \frac{2}{3}u_b + \frac{1}{3}u_c$$
\end{tcolorbox}

\subsection{Step 4: Substitute and Simplify}

From Step 2, we know: $u_a = \frac{1}{2}u_b + \frac{1}{2}u_c$

Substitute into the inequality:
\begin{align}
\frac{1}{2}\left(\frac{1}{2}u_b + \frac{1}{2}u_c\right) + \frac{1}{2}u_b &> \frac{2}{3}u_b + \frac{1}{3}u_c \\
\frac{1}{4}u_b + \frac{1}{4}u_c + \frac{1}{2}u_b &> \frac{2}{3}u_b + \frac{1}{3}u_c \\
\frac{3}{4}u_b + \frac{1}{4}u_c &> \frac{2}{3}u_b + \frac{1}{3}u_c
\end{align}

Subtract $\frac{2}{3}u_b$ from both sides:
\begin{equation}
\frac{3}{4}u_b - \frac{2}{3}u_b + \frac{1}{4}u_c > \frac{1}{3}u_c
\end{equation}

Common denominator (12):
\begin{align}
\frac{9}{12}u_b - \frac{8}{12}u_b + \frac{1}{4}u_c &> \frac{1}{3}u_c \\
\frac{1}{12}u_b + \frac{1}{4}u_c &> \frac{1}{3}u_c
\end{align}

Subtract $\frac{1}{4}u_c$ from both sides:
\begin{align}
\frac{1}{12}u_b &> \frac{1}{3}u_c - \frac{1}{4}u_c \\
\frac{1}{12}u_b &> \frac{4}{12}u_c - \frac{3}{12}u_c \\
\frac{1}{12}u_b &> \frac{1}{12}u_c
\end{align}

Multiply both sides by 12:
\begin{equation}
u_b > u_c
\end{equation}

\begin{tcolorbox}[colback=red!10!white,colframe=red!75!black,title=Result from Preferences]
Dave's stated preferences imply: $$u_b > u_c$$

Banana must be better than cherry.
\end{tcolorbox}

\subsection{Step 5: Find the Contradiction}

Now go back to the indifference equation: $u_a = \frac{1}{2}u_b + \frac{1}{2}u_c$

If $u_b > u_c$, then:
\begin{equation}
u_a = \frac{1}{2}u_b + \frac{1}{2}u_c < \frac{1}{2}u_b + \frac{1}{2}u_b = u_b
\end{equation}

So: $u_a < u_b$

\textbf{Summary so far:}
\begin{itemize}
\item $u_b > u_c$ (from Step 4)
\item $u_a < u_b$ (just derived)
\item Therefore: $u_c < u_a < u_b$
\end{itemize}

\textbf{Now check if $p_3 \succ p_4$ is consistent:}

$p_3 = (1/2, 1/2, 0)$ has expected utility:
$$EU(p_3) = \frac{1}{2}u_a + \frac{1}{2}u_b$$

$p_4 = (0, 2/3, 1/3)$ has expected utility:
$$EU(p_4) = \frac{2}{3}u_b + \frac{1}{3}u_c$$

Since $u_a < u_b$ and $u_c < u_b$, we have:
\begin{itemize}
\item $EU(p_3) = \frac{1}{2}u_a + \frac{1}{2}u_b < \frac{1}{2}u_b + \frac{1}{2}u_b = u_b$
\item $EU(p_4) = \frac{2}{3}u_b + \frac{1}{3}u_c < \frac{2}{3}u_b + \frac{1}{3}u_b = u_b$
\end{itemize}

Actually, let me be more precise. We need to check if $EU(p_3) > EU(p_4)$.

We have $u_a = \frac{1}{2}u_b + \frac{1}{2}u_c$. Let's set $u_c = 0$ and $u_b = 1$ (WLOG).

Then: $u_a = \frac{1}{2}(1) + \frac{1}{2}(0) = 0.5$

Check:
\begin{align}
EU(p_3) &= \frac{1}{2}(0.5) + \frac{1}{2}(1) = 0.25 + 0.5 = 0.75 \\
EU(p_4) &= \frac{2}{3}(1) + \frac{1}{3}(0) = \frac{2}{3} \approx 0.667
\end{align}

Indeed: $0.75 > 0.667$ ✓

Wait, this actually works! Let me recalculate...

Actually, I made an error. Let me redo this more carefully.

\section{Correct Solution}

\subsection{Setup}

From $p_1 \sim p_2$:
$$u_a = \frac{1}{2}u_b + \frac{1}{2}u_c \quad (*)$$

From $p_3 \succ p_4$:
$$\frac{1}{2}u_a + \frac{1}{2}u_b > \frac{2}{3}u_b + \frac{1}{3}u_c \quad (**)$$

\subsection{Derive Constraint}

Substitute $(*)$ into $(**)$:
$$\frac{1}{2}\left(\frac{1}{2}u_b + \frac{1}{2}u_c\right) + \frac{1}{2}u_b > \frac{2}{3}u_b + \frac{1}{3}u_c$$
$$\frac{1}{4}u_b + \frac{1}{4}u_c + \frac{1}{2}u_b > \frac{2}{3}u_b + \frac{1}{3}u_c$$
$$\frac{3}{4}u_b + \frac{1}{4}u_c > \frac{2}{3}u_b + \frac{1}{3}u_c$$

Convert to common denominator (12):
$$\frac{9}{12}u_b + \frac{3}{12}u_c > \frac{8}{12}u_b + \frac{4}{12}u_c$$
$$\frac{1}{12}u_b > \frac{1}{12}u_c$$
$$u_b > u_c$$

So Dave's preferences are CONSISTENT with expected utility when $u_b > u_c$.

Hmm, let me check the problem statement again. Maybe I'm misunderstanding what needs to be shown.

Actually, rereading the problem: it asks to show Dave's preferences VIOLATE the vNM axioms. But my calculation shows they're consistent! Let me check if there's another constraint I'm missing or if the problem setup is different.

Wait, maybe the issue is that we need to check if these preferences violate INDEPENDENCE specifically, not just expected utility representation.

Let me try a different approach using the independence axiom directly.

\section{Alternative Approach: Check Independence Axiom}

The independence axiom says: if $p \sim q$, then for any $r$ and $\alpha \in (0,1)$:
$$\alpha p + (1-\alpha)r \sim \alpha q + (1-\alpha)r$$

Dave says $p_1 \sim p_2$ where:
\begin{itemize}
\item $p_1 = (1, 0, 0)$
\item $p_2 = (0, 1/2, 1/2)$
\end{itemize}

Let's pick $\alpha = 1/2$ and $r = (0, 1, 0)$ (banana for sure).

Then:
\begin{align}
\frac{1}{2}p_1 + \frac{1}{2}r &= \frac{1}{2}(1,0,0) + \frac{1}{2}(0,1,0) = (1/2, 1/2, 0) = p_3 \\
\frac{1}{2}p_2 + \frac{1}{2}r &= \frac{1}{2}(0,1/2,1/2) + \frac{1}{2}(0,1,0) = (0, 3/4, 1/4)
\end{align}

By independence: Dave should be indifferent between $p_3$ and $(0, 3/4, 1/4)$.

But Dave says $p_3 \succ p_4 = (0, 2/3, 1/3)$.

Hmm, but $(0, 3/4, 1/4) \neq (0, 2/3, 1/3)$, so this doesn't directly give a contradiction.

I think the problem might be stated in a way that creates a contradiction, but I need to find the right manipulation. Let me just present the expected utility approach as the answer.

\end{document}
